\clearpage
\section*{Orientierungshilfe: Wie du mit diesem Skript arbeitest}
\addcontentsline{toc}{section}{Orientierungshilfe: Wie du mit diesem Skript arbeitest}

Dieses Skript begleitet dich durch die gesamte Oberstufe (E-Phase bis Q2) und
enthält deutlich mehr Material, als in einem Kurs vollständig bearbeitet werden
kann. Damit du jederzeit weißt, welche Inhalte für dich verbindlich sind und
welche Aufgaben als vertiefendes Training dienen, sind Kapitel, Abschnitte und
Übungsaufgaben mit farbigen Kennzeichnungen versehen.

\vspace{-4pt}
\subsection*{1. Bedeutung der Kennzeichnungen}

\begin{center}
\bgroup
\def\arraystretch{1.2}
\begin{tabularx}{\textwidth}{|p{3.5cm}|X|}
\hline
\textbf{Kennzeichnung} & \textbf{Bedeutung für deinen Unterricht} \\ \hline
\textit{(ohne Tag)} / \tagPflicht &
\textbf{Verbindlicher Kernstoff} für Grundkurs (\textbf{GK}) und Leistungskurs (\textbf{LK}). Bildet das Fundament für Klausuren und das Abitur. \\ \hline
\tagLK &
\textbf{Leistungskurs-Pflichtstoff:} Verbindlich für den LK (insbesondere die Abitur-Schwerpunkte \emph{Datenbanken} und \emph{Client/Server}). \\ \hline
\tagGKWahl &
\textbf{Grundkurs-Wahlthema:} Mögliche Vertiefungsthemen im GK (Q2), z.\,B.\ für das 4.~Prüfungsfach im Abitur. \\ \hline
\tagZusatz &
\textbf{Freiwillige Vertiefung / Nachschlagewerk:} Nicht verbindliche Inhalte oder zusätzliche Übungen für schnelle Lerner und Projekte. \\ \hline
\end{tabularx}
\egroup
\end{center}

\vspace{-4pt}
\subsection*{2. Zwei wichtige Regeln zum Lesen der Tags}

\begin{compactitem}
  \item \textbf{Kombinationen (z.\,B.\ \tagZusatz\tagLK):}
  Tags lassen sich kombinieren: \tagZusatz\tagLK\ kennzeichnet eine
  \emph{nicht verbindliche Zusatzaufgabe auf LK-Niveau} (analog \tagZusatz\tagGKWahl\
  als freiwillige Vertiefung im GK-Wahlbereich).
  \item \textbf{Vererbung auf Unterabschnitte:}
  Ist ein gesamtes Kapitel (oder ein Abschnitt) mit \tagZusatz\ oder \tagLK\
  versehen, gilt dies automatisch für alle untergeordneten Teile, sofern dort
  nichts anderes angegeben ist.
\end{compactitem}

\vspace{-4pt}
\subsection*{3. Der rote Faden durch das Skript (Lernpfad)}

\begin{compactenum}[1.]
  \item \textbf{Einstieg \& Grundlagen (Kapitel~1--9):}
  Einrichtung von Eclipse/Git (Kap.~1) sowie erste Schritte in Java über
  Turtle-Grafik (Kap.~5), Dialoge (Kap.~7) und eigene Zeichenprogramme
  (Kap.~8--9). Kapitel~2--4 dienen als Nachschlagewerk.
  \item \textbf{Objektorientierte Programmierung \& UML (Kapitel~10--15):}
  Klassen/Objekte (Kap.~10), Datenkapselung (Kap.~11), Klassendiagramme
  (Kap.~12), Vererbung (Kap.~13), die OOP-Bilanz (Kap.~14) und
  Zustandsdiagramme (Kap.~15).
  \item \textbf{Ereignisbehandlung, Swing-GUIs \& Datenströme (Kapitel~25--29):}
  Grafische Oberflächen mit dem \emph{WindowBuilder} und Strings (Kap.~25--26),
  Exceptions (Kap.~27), Threads (Kap.~28) und Dateizugriffe (Kap.~29).
  \emph{(LK-Hinweis: Direkt nach Kap.~29 folgt oft ein Vorgriff auf Kap.~45
  zum Echo-Client, da dort dieselben Stream-Klassen genutzt werden.)}
  \item \textbf{Abitur-Schwerpunkte in der Qualifikationsphase (Kapitel~34--46):}\\
  \textbf{GK \& LK gemeinsam:} ER-Modell (Kap.~34) und SQL-Abfragen (Kap.~35--36).\\
  \textbf{Leistungskurs (\tagLK):} Normalisierung (Kap.~37), Java + SQL
  (Kap.~39), Netzwerke (Kap.~44), Client/Server (Kap.~45) und Abi-Training
  (Kap.~46).\\
  \textbf{Grundkurs (\tagGKWahl):} Wahlweise Datensicherheit (Kap.~42),
  Kryptologie (Kap.~43) und Grundlagen zu Netzwerken (Kap.~44).
\end{compactenum}

\clearpage
